\section{Related Work}
\label{sec:related}

As the number of CPU cores has scaled up in recent years, the traditional design of utilizing a single log buffer and a single log flusher has increasingly become a critical performance bottleneck, primarily due to severe contention over the log stream. To mitigate this issue, extensive research has been conducted on parallel logging or decentralized logging, where multiple cores record and append log records to persistent storage concurrently.

Pioneering efforts in parallel logging include systems such as Ather \cite{} and passive group commit \cite{}. These approaches allocate isolated log buffers (partitioned logs) per core or per socket, thereby completely eliminating latch contention during log insertion. 
However, a naive partitioning of logs introduces a fundamental trade-off: it complicates the precise reconstruction of global dependencies—such as the exact commit order of transactions—during crash recovery. To manage these dependencies efficiently, SiloR \cite{zheng2014silor} introduced a cooperative, epoch-based mechanism that flushes the log records of all cores collectively at fixed time intervals (epochs), ensuring consistency at the granularity of an epoch. Meanwhile, Taurus \cite{xia2021taurus} embeds logical timestamps into log records to capture data dependencies among transactions. By dynamically reconstructing a dependency graph during the recovery phase, Taurus successfully achieves both out-of-order parallel log writes and accurate recovery.

Furthermore, the emergence of non-volatile RAM (NVRAM) has brought a paradigm shift to the design of parallel logging protocols. Because NVRAM offers byte-addressability and low latency, it challenges the fundamental premise of traditional Write-Ahead Logging (WAL), which relies on translating in-memory data modifications into a serialized, disk-oriented log stream.

For instance, Write-Behind Logging (WBL) \cite{wbl} leverages the persistence of NVRAM by flushing the modified data itself directly to the persistent layer before a transaction commits. By recording only minimal commit metadata in the log, WBL drastically reduces the overhead of generating Redo logs. Other studies, such as Rethinking Logging for NVRAM \cite{rethinking} or similar fine-grained approaches, propose parallel, fine-grained checkpointing on NVRAM, which either eliminates the need for operational logs entirely or significantly restricts their role. Additionally, systems like XXX \cite{XXX} allocate independent, per-core log spaces directly within NVRAM, enabling each core to perform synchronous writes without cross-core coordination, thereby achieving parallel logging throughput that approaches the physical limits of the hardware.

In summary, modern logging mechanisms continue to evolve toward a harmonious integration of multi-core parallelism and the byte-addressable persistence provided by NVRAM.